\section{有效性威胁与局限}
\label{sec:threats}

\paragraph{构造有效性（construct validity）}
五个夹具覆盖上下界、空循环、正常循环、递归和短路副作用，但输出一致不构成全输入形式证明。退出状态与精确标准输出是合理的外部判据，仍可能漏掉不影响这些输入的内存错误。可进一步用解释器、未定义行为检查器或等价性验证扩展证据，但不能把它们倒推为本实验已经完成。

\paragraph{内部有效性（internal validity）}
源码路径由 Clang 以 C 语法接收 SysY 子集，未证明任意 SysY 编译器都作相同处理。通过限制程序到两者公共语义、固定 Clang 18.1.3、目标三元组和 \texttt{rv64gc/lp64d}，可降低工具和语言差异的混杂。\texttt{-O0}/\texttt{-O2} 对照只支持 IR/汇编结构结论；未做计时，故不支持性能因果判断。

\paragraph{外部有效性（external validity）}
一个八元素整数程序不足以代表浮点、并发、大型链接或复杂优化。QEMU 用户态验证的是 RV64 Linux 程序行为，不等于真实微架构性能，也不适用于裸机 ABI。预编译运行时的兼容性判断针对所审计归档；若其来源或构建参数变化，应重新检查符号与 libc 家族。

\paragraph{进阶边界}
AscendNPU IR 结论来自静态文本和官方文档。结构检查不会验证方言注册、转换合法性、设备对象生成、数值结果或性能；缺少专用工具链和硬件时，停在这个边界比把文档预期输出冒充观测更可靠。
